% Options for packages loaded elsewhere
\PassOptionsToPackage{unicode}{hyperref}
\PassOptionsToPackage{hyphens}{url}
\documentclass[
]{article}
\usepackage{xcolor}
\usepackage[margin=1in]{geometry}
\usepackage{amsmath,amssymb}
\setcounter{secnumdepth}{-\maxdimen} % remove section numbering
\usepackage{iftex}
\ifPDFTeX
  \usepackage[T1]{fontenc}
  \usepackage[utf8]{inputenc}
  \usepackage{textcomp} % provide euro and other symbols
\else % if luatex or xetex
  \usepackage{unicode-math} % this also loads fontspec
  \defaultfontfeatures{Scale=MatchLowercase}
  \defaultfontfeatures[\rmfamily]{Ligatures=TeX,Scale=1}
\fi
\usepackage{lmodern}
\ifPDFTeX\else
  % xetex/luatex font selection
\fi
% Use upquote if available, for straight quotes in verbatim environments
\IfFileExists{upquote.sty}{\usepackage{upquote}}{}
\IfFileExists{microtype.sty}{% use microtype if available
  \usepackage[]{microtype}
  \UseMicrotypeSet[protrusion]{basicmath} % disable protrusion for tt fonts
}{}
\makeatletter
\@ifundefined{KOMAClassName}{% if non-KOMA class
  \IfFileExists{parskip.sty}{%
    \usepackage{parskip}
  }{% else
    \setlength{\parindent}{0pt}
    \setlength{\parskip}{6pt plus 2pt minus 1pt}}
}{% if KOMA class
  \KOMAoptions{parskip=half}}
\makeatother
\usepackage{color}
\usepackage{fancyvrb}
\newcommand{\VerbBar}{|}
\newcommand{\VERB}{\Verb[commandchars=\\\{\}]}
\DefineVerbatimEnvironment{Highlighting}{Verbatim}{commandchars=\\\{\}}
% Add ',fontsize=\small' for more characters per line
\usepackage{framed}
\definecolor{shadecolor}{RGB}{248,248,248}
\newenvironment{Shaded}{\begin{snugshade}}{\end{snugshade}}
\newcommand{\AlertTok}[1]{\textcolor[rgb]{0.94,0.16,0.16}{#1}}
\newcommand{\AnnotationTok}[1]{\textcolor[rgb]{0.56,0.35,0.01}{\textbf{\textit{#1}}}}
\newcommand{\AttributeTok}[1]{\textcolor[rgb]{0.13,0.29,0.53}{#1}}
\newcommand{\BaseNTok}[1]{\textcolor[rgb]{0.00,0.00,0.81}{#1}}
\newcommand{\BuiltInTok}[1]{#1}
\newcommand{\CharTok}[1]{\textcolor[rgb]{0.31,0.60,0.02}{#1}}
\newcommand{\CommentTok}[1]{\textcolor[rgb]{0.56,0.35,0.01}{\textit{#1}}}
\newcommand{\CommentVarTok}[1]{\textcolor[rgb]{0.56,0.35,0.01}{\textbf{\textit{#1}}}}
\newcommand{\ConstantTok}[1]{\textcolor[rgb]{0.56,0.35,0.01}{#1}}
\newcommand{\ControlFlowTok}[1]{\textcolor[rgb]{0.13,0.29,0.53}{\textbf{#1}}}
\newcommand{\DataTypeTok}[1]{\textcolor[rgb]{0.13,0.29,0.53}{#1}}
\newcommand{\DecValTok}[1]{\textcolor[rgb]{0.00,0.00,0.81}{#1}}
\newcommand{\DocumentationTok}[1]{\textcolor[rgb]{0.56,0.35,0.01}{\textbf{\textit{#1}}}}
\newcommand{\ErrorTok}[1]{\textcolor[rgb]{0.64,0.00,0.00}{\textbf{#1}}}
\newcommand{\ExtensionTok}[1]{#1}
\newcommand{\FloatTok}[1]{\textcolor[rgb]{0.00,0.00,0.81}{#1}}
\newcommand{\FunctionTok}[1]{\textcolor[rgb]{0.13,0.29,0.53}{\textbf{#1}}}
\newcommand{\ImportTok}[1]{#1}
\newcommand{\InformationTok}[1]{\textcolor[rgb]{0.56,0.35,0.01}{\textbf{\textit{#1}}}}
\newcommand{\KeywordTok}[1]{\textcolor[rgb]{0.13,0.29,0.53}{\textbf{#1}}}
\newcommand{\NormalTok}[1]{#1}
\newcommand{\OperatorTok}[1]{\textcolor[rgb]{0.81,0.36,0.00}{\textbf{#1}}}
\newcommand{\OtherTok}[1]{\textcolor[rgb]{0.56,0.35,0.01}{#1}}
\newcommand{\PreprocessorTok}[1]{\textcolor[rgb]{0.56,0.35,0.01}{\textit{#1}}}
\newcommand{\RegionMarkerTok}[1]{#1}
\newcommand{\SpecialCharTok}[1]{\textcolor[rgb]{0.81,0.36,0.00}{\textbf{#1}}}
\newcommand{\SpecialStringTok}[1]{\textcolor[rgb]{0.31,0.60,0.02}{#1}}
\newcommand{\StringTok}[1]{\textcolor[rgb]{0.31,0.60,0.02}{#1}}
\newcommand{\VariableTok}[1]{\textcolor[rgb]{0.00,0.00,0.00}{#1}}
\newcommand{\VerbatimStringTok}[1]{\textcolor[rgb]{0.31,0.60,0.02}{#1}}
\newcommand{\WarningTok}[1]{\textcolor[rgb]{0.56,0.35,0.01}{\textbf{\textit{#1}}}}
\usepackage{graphicx}
\makeatletter
\newsavebox\pandoc@box
\newcommand*\pandocbounded[1]{% scales image to fit in text height/width
  \sbox\pandoc@box{#1}%
  \Gscale@div\@tempa{\textheight}{\dimexpr\ht\pandoc@box+\dp\pandoc@box\relax}%
  \Gscale@div\@tempb{\linewidth}{\wd\pandoc@box}%
  \ifdim\@tempb\p@<\@tempa\p@\let\@tempa\@tempb\fi% select the smaller of both
  \ifdim\@tempa\p@<\p@\scalebox{\@tempa}{\usebox\pandoc@box}%
  \else\usebox{\pandoc@box}%
  \fi%
}
% Set default figure placement to htbp
\def\fps@figure{htbp}
\makeatother
\setlength{\emergencystretch}{3em} % prevent overfull lines
\providecommand{\tightlist}{%
  \setlength{\itemsep}{0pt}\setlength{\parskip}{0pt}}
\usepackage{bookmark}
\IfFileExists{xurl.sty}{\usepackage{xurl}}{} % add URL line breaks if available
\urlstyle{same}
\hypersetup{
  pdftitle={Assignment\_2\_Statistics\_II},
  pdfauthor={Maik David Klaiber (22-482-426); Jonas Weber (24-621-484); Yann Rochaix (24-622-276)},
  hidelinks,
  pdfcreator={LaTeX via pandoc}}

\title{Assignment\_2\_Statistics\_II}
\author{Maik David Klaiber (22-482-426) \and Jonas Weber
(24-621-484) \and Yann Rochaix (24-622-276)}
\date{2026-03-14}

\begin{document}
\maketitle

\subsection{Exercise 1 - Data
Preparation}\label{exercise-1---data-preparation}

We imported the datasets using \texttt{read\_delim()} and
\texttt{read\_csv()} from the \texttt{readr} package to handle different
separators. We then used \texttt{dfSummary()} from the
\texttt{summarytools} package to automatically check the data. The
difference in row counts gives us a first hint, that there might be a
mismatch in the number of cantons. This summary also shows us, that
there are no missing values.

\begin{Shaded}
\begin{Highlighting}[]
\NormalTok{cantons}\OtherTok{\textless{}{-}} \FunctionTok{read\_delim}\NormalTok{(}\StringTok{"ps2\_cantons\_cleaned.csv"}\NormalTok{, }\AttributeTok{delim =} \StringTok{";"}\NormalTok{)}
\NormalTok{surnames}\OtherTok{\textless{}{-}} \FunctionTok{read\_csv}\NormalTok{(}\StringTok{"ps2\_cantons\_surnames.csv"}\NormalTok{)}
\FunctionTok{dfSummary}\NormalTok{(surnames)}
\FunctionTok{dfSummary}\NormalTok{(cantons)}
\end{Highlighting}
\end{Shaded}

We used \texttt{mutate()} to consolidate the individual proportions into
a single surname variable and then applied \texttt{select()} to remove
the original columns, streamlining the tibble for a cleaner merge. Then,
to show the rows that have no matching canton name in the canton
dataset, we use \texttt{antijoin} to reveal them.

\begin{Shaded}
\begin{Highlighting}[]
\NormalTok{surnames }\OtherTok{\textless{}{-}}\NormalTok{ surnames }\SpecialCharTok{\%\textgreater{}\%}
  \FunctionTok{mutate}\NormalTok{(}\AttributeTok{surname =}\NormalTok{ share\_mueller }\SpecialCharTok{+}\NormalTok{ share\_maier) }\SpecialCharTok{\%\textgreater{}\%}
  \FunctionTok{select}\NormalTok{(}\SpecialCharTok{{-}}\NormalTok{share\_mueller, }\SpecialCharTok{{-}}\NormalTok{share\_maier)}

\NormalTok{surnames }\SpecialCharTok{\%\textgreater{}\%} 
  \FunctionTok{anti\_join}\NormalTok{(cantons, }\AttributeTok{by =} \StringTok{"canton"}\NormalTok{)}
\end{Highlighting}
\end{Shaded}

\begin{verbatim}
## # A tibble: 2 x 2
##   canton        surname
##   <chr>           <dbl>
## 1 Sri Lanka    0.000016
## 2 Sankt Gallen 0.0726
\end{verbatim}

In this step, we harmonized the naming convention by recoding ``Sankt
Gallen'' to ``St.~Gallen'' and filtered out the ``Sri Lanka'' entry to
ensure the surnames dataset exclusively contains valid Swiss cantons for
the subsequent join. After performing the merge, we used
\texttt{select()} to retain only the essential variables for analysis:
GDP, population, immigration, and the newly calculated surname shares.
We then verified the integrity of the join by confirming there are zero
missing values and that all 26 distinct Swiss cantons are correctly
represented in the final data frame.

\begin{Shaded}
\begin{Highlighting}[]
\NormalTok{surnames }\OtherTok{\textless{}{-}}\NormalTok{ surnames }\SpecialCharTok{\%\textgreater{}\%}
  \FunctionTok{mutate}\NormalTok{(}\AttributeTok{canton =} \FunctionTok{recode}\NormalTok{(canton, }\StringTok{"Sankt Gallen"} \OtherTok{=} \StringTok{"St. Gallen"}\NormalTok{)) }\SpecialCharTok{\%\textgreater{}\%}
  \FunctionTok{filter}\NormalTok{(canton }\SpecialCharTok{!=} \StringTok{"Sri Lanka"}\NormalTok{)}

\NormalTok{df }\OtherTok{\textless{}{-}}\NormalTok{ cantons }\SpecialCharTok{\%\textgreater{}\%} \CommentTok{\#Merge the surname dataset into the canton dataset}
  \FunctionTok{left\_join}\NormalTok{(surnames, }\AttributeTok{by =} \StringTok{"canton"}\NormalTok{) }\SpecialCharTok{\%\textgreater{}\%}
  \FunctionTok{select}\NormalTok{(gdp, population, immigration, surname)}

\FunctionTok{sum}\NormalTok{(}\FunctionTok{is.na}\NormalTok{(df}\SpecialCharTok{$}\NormalTok{surname))          }
\end{Highlighting}
\end{Shaded}

\begin{verbatim}
## [1] 0
\end{verbatim}

\begin{Shaded}
\begin{Highlighting}[]
\FunctionTok{n\_distinct}\NormalTok{(cantons}\SpecialCharTok{$}\NormalTok{canton)}
\end{Highlighting}
\end{Shaded}

\begin{verbatim}
## [1] 26
\end{verbatim}

\subsection{Exercise 2 - Estimation of an univariate
model}\label{exercise-2---estimation-of-an-univariate-model}

The specific regression equation for this problem is:
\[gdp_i = \beta_0 + \beta_1 \cdot surname_i + u_i\]

\begin{Shaded}
\begin{Highlighting}[]
\NormalTok{model }\OtherTok{\textless{}{-}} \FunctionTok{lm}\NormalTok{(gdp }\SpecialCharTok{\textasciitilde{}}\NormalTok{ surname, }\AttributeTok{data =}\NormalTok{ df)}
\FunctionTok{summary}\NormalTok{(model)}
\end{Highlighting}
\end{Shaded}

\begin{verbatim}
## 
## Call:
## lm(formula = gdp ~ surname, data = df)
## 
## Residuals:
##    Min     1Q Median     3Q    Max 
## -43665 -10345  -7762   6033  42831 
## 
## Coefficients:
##             Estimate Std. Error t value Pr(>|t|)    
## (Intercept)    10077       4856   2.075   0.0489 *  
## surname       449550      68427   6.570 8.54e-07 ***
## ---
## Signif. codes:  0 '***' 0.001 '**' 0.01 '*' 0.05 '.' 0.1 ' ' 1
## 
## Residual standard error: 20130 on 24 degrees of freedom
## Multiple R-squared:  0.6427, Adjusted R-squared:  0.6278 
## F-statistic: 43.16 on 1 and 24 DF,  p-value: 8.54e-07
\end{verbatim}

\subsubsection{Results}\label{results}

\textbf{Intercept \(\hat{\beta}_0 = 10,077.28\):} A canton with a
\texttt{surname} share of zero is predicted to have a GDP of CHF 10'077
million. This is not meaningful in practice since all cantons have some
share of these surnames.

\textbf{Slope \(\hat{\beta}_1 = 449,550\):} For each one percentage
point increase in the combined Mueller/Maier surname share, GDP is
predicted to increase by CHF 449'550 million. This reflects that larger
cantons naturally have both more people with common surnames and higher
total GDP.

\begin{Shaded}
\begin{Highlighting}[]
\NormalTok{avg\_surname\_weighted }\OtherTok{\textless{}{-}} \FunctionTok{sum}\NormalTok{(df}\SpecialCharTok{$}\NormalTok{surname }\SpecialCharTok{*}\NormalTok{ df}\SpecialCharTok{$}\NormalTok{population) }\SpecialCharTok{/} \FunctionTok{sum}\NormalTok{(df}\SpecialCharTok{$}\NormalTok{population)}
\NormalTok{predicted\_gdp }\OtherTok{\textless{}{-}} \FunctionTok{predict}\NormalTok{(model, }\AttributeTok{newdata =} \FunctionTok{data.frame}\NormalTok{(}\AttributeTok{surname =}\NormalTok{ avg\_surname\_weighted))}
\FunctionTok{print}\NormalTok{(}\FunctionTok{as.numeric}\NormalTok{(predicted\_gdp))}
\end{Highlighting}
\end{Shaded}

\begin{verbatim}
## [1] 51512.57
\end{verbatim}

For each additional percentage point increase in the Mueller/Maier
surname share beyond this average, GDP is predicted to increase by CHF
4'495,50 million.

\subsection{Exercise 3 - Comparison of two univariate
models}\label{exercise-3---comparison-of-two-univariate-models}

We estimate the regression of \(gdp\) on \(\log(\text{immigration})\)
and compute the corresponding \(R^2\).

\begin{Shaded}
\begin{Highlighting}[]
\CommentTok{\# Estimate regression of GDP on the log of immigration}
\NormalTok{model\_log }\OtherTok{\textless{}{-}} \FunctionTok{lm}\NormalTok{(gdp }\SpecialCharTok{\textasciitilde{}} \FunctionTok{log}\NormalTok{(immigration), }\AttributeTok{data =}\NormalTok{ cantons)}
\FunctionTok{summary}\NormalTok{(model\_log)}
\end{Highlighting}
\end{Shaded}

\begin{verbatim}
## 
## Call:
## lm(formula = gdp ~ log(immigration), data = cantons)
## 
## Residuals:
##    Min     1Q Median     3Q    Max 
## -25926 -11024  -5529   5363  79253 
## 
## Coefficients:
##                  Estimate Std. Error t value Pr(>|t|)    
## (Intercept)          8147       5353   1.522    0.141    
## log(immigration)    18564       3061   6.065  2.9e-06 ***
## ---
## Signif. codes:  0 '***' 0.001 '**' 0.01 '*' 0.05 '.' 0.1 ' ' 1
## 
## Residual standard error: 21160 on 24 degrees of freedom
## Multiple R-squared:  0.6052, Adjusted R-squared:  0.5887 
## F-statistic: 36.79 on 1 and 24 DF,  p-value: 2.903e-06
\end{verbatim}

\begin{Shaded}
\begin{Highlighting}[]
\CommentTok{\# Define R² from the population model estimated in Exercise 2}
\NormalTok{r\_squared\_surname }\OtherTok{\textless{}{-}} \FloatTok{0.6427}

\CommentTok{\# Create a small table comparing the R² values of both models}
\NormalTok{r2\_table }\OtherTok{\textless{}{-}} \FunctionTok{tibble}\NormalTok{(}
  \AttributeTok{model =} \FunctionTok{c}\NormalTok{(}\StringTok{"surname (exercise 2)"}\NormalTok{, }\StringTok{"log(immigration) model"}\NormalTok{),}
  \AttributeTok{R\_squared =} \FunctionTok{c}\NormalTok{(r\_squared\_surname, }\FunctionTok{summary}\NormalTok{(model\_log)}\SpecialCharTok{$}\NormalTok{r.squared)}
\NormalTok{)}

\NormalTok{r2\_table}
\end{Highlighting}
\end{Shaded}

\begin{verbatim}
## # A tibble: 2 x 2
##   model                  R_squared
##   <chr>                      <dbl>
## 1 surname (exercise 2)       0.643
## 2 log(immigration) model     0.605
\end{verbatim}

What does \(R^2\) measure here intuitively? \(R^2\) measures the
proportion of the variation in \(y\) that is explained by the regression
model. A higher \(R^2\) indicates greater explanatory power of the
model.

What amount of variation in \texttt{gdp} is explained? The surname model
explains about 64.3\% of the variation in \(gdp\), while the
log(immigration) model explains about 60.5\%. Thus, the surname model
has higher explanatory power.

Unbiasedness of the Estimators for Surname and Immigration The OLS
estimators are unbiased if the assumption \(E(u \mid X) = 0\) holds
(SLR.4), meaning that the error term has an average value of zero for
any given value of \(X\). This requires that the explanatory variable is
not systematically related to omitted factors that affect \(gdp\).

In the surname model, this assumption is likely violated because surname
shares may be correlated with underlying regional or demographic
characteristics that also influence economic performance. If such
relevant factors are omitted from the regression, the estimator for the
surname coefficient may be biased due to omitted variable bias.

A similar issue may arise in the immigration model. For example,
\texttt{education\ levels} could be an omitted variable. Higher
education levels can increase productivity and therefore \(gdp\). At the
same time, economically attractive cantons with better labor market
opportunities may attract more immigrants. If education is correlated
with immigration and affects \(gdp\), the immigration coefficient may
capture part of this effect, leading to biased OLS estimates.

Relationship between unbiasedness of a coefficient estimator and \(R^2\)
Unbiasedness of an estimator depends on the assumption
\(E(u \mid X) = 0\). In contrast, \(R^2\) measures how much of the
variation in the dependent variable is explained by the model.
Therefore, unbiasedness and predictive power are not directly related
and a model can have a high \(R^2\) while the coefficient estimates are
still biased.

\subsection{Exercise 4 - Univariate model without a
constant}\label{exercise-4---univariate-model-without-a-constant}

We state the formula without the constant \(\beta_0\) as \[
gdp_i = \beta_1 \, surname_i + u_i
\]

\begin{Shaded}
\begin{Highlighting}[]
\CommentTok{\# Estimate the regression of gdp on surname without an intercept ({-}1)}
\FunctionTok{lm}\NormalTok{(gdp }\SpecialCharTok{\textasciitilde{}}\NormalTok{ surname }\SpecialCharTok{{-}} \DecValTok{1}\NormalTok{, }\AttributeTok{data =}\NormalTok{ df)}
\end{Highlighting}
\end{Shaded}

\begin{verbatim}
## 
## Call:
## lm(formula = gdp ~ surname - 1, data = df)
## 
## Coefficients:
## surname  
##  532247
\end{verbatim}

Why did the coefficient on surname change compared to subexercise 2? The
coefficient changes because the regression is estimated without an
intercept, forcing the regression line to pass through the origin. In
this case, the OLS estimator simplifies from
\(\hat{\beta}_1 = \frac{\sum (x_i-\bar{x})(y_i-\bar{y})}{\sum (x_i-\bar{x})^2}\)
to \(\hat{\beta}_1 = \frac{\sum x_i y_i}{\sum x_i^2}\), which changes
the estimated slope (from 449,550 to 532,247).

\subsection{Exercise 5 - Code your own OLS
estimator}\label{exercise-5---code-your-own-ols-estimator}

We estimate the parameters of the univariate regression model

\[
y_i = \beta_0 + \beta_1 x_i + u_i .
\]

Using the OLS formulas, the slope coefficient is given by

\[
\hat{\beta}_1 =
\frac{\sum_{i=1}^{N}(x_i-\bar{x})(y_i-\bar{y})}
{\sum_{i=1}^{N}(x_i-\bar{x})^2}.
\]

The intercept is then calculated as

\[
\hat{\beta}_0 = \bar{y} - \hat{\beta}_1 \bar{x}.
\]

In our application, we regress \texttt{gdp} (as Y) on \texttt{surname}
(as X), including an intercept. Thus, the vectors \texttt{gdp} and
\texttt{surname} are used as inputs to compute the OLS estimators using
the formulas above.

\begin{Shaded}
\begin{Highlighting}[]
\NormalTok{estimate\_ols }\OtherTok{\textless{}{-}} \ControlFlowTok{function}\NormalTok{(y, x) \{}
  
  \CommentTok{\# Compute the OLS slope coefficient beta\_1 using the analytical formula above}
\NormalTok{  beta\_1 }\OtherTok{\textless{}{-}} \FunctionTok{sum}\NormalTok{((x }\SpecialCharTok{{-}} \FunctionTok{mean}\NormalTok{(x)) }\SpecialCharTok{*}\NormalTok{ (y }\SpecialCharTok{{-}} \FunctionTok{mean}\NormalTok{(y))) }\SpecialCharTok{/} \FunctionTok{sum}\NormalTok{((x }\SpecialCharTok{{-}} \FunctionTok{mean}\NormalTok{(x))}\SpecialCharTok{\^{}}\DecValTok{2}\NormalTok{)}
  
  \CommentTok{\# Compute the intercept beta\_0}
\NormalTok{  beta\_0 }\OtherTok{\textless{}{-}} \FunctionTok{mean}\NormalTok{(y) }\SpecialCharTok{{-}}\NormalTok{ beta\_1 }\SpecialCharTok{*} \FunctionTok{mean}\NormalTok{(x)}
  
  \CommentTok{\# Return the estimated coefficients as a named vector}
  \FunctionTok{return}\NormalTok{(}\FunctionTok{c}\NormalTok{(}\AttributeTok{beta\_0 =}\NormalTok{ beta\_0, }\AttributeTok{beta\_1 =}\NormalTok{ beta\_1))}
\NormalTok{\}}

\FunctionTok{estimate\_ols}\NormalTok{(df}\SpecialCharTok{$}\NormalTok{gdp, df}\SpecialCharTok{$}\NormalTok{surname)}
\end{Highlighting}
\end{Shaded}

\begin{verbatim}
##    beta_0    beta_1 
##  10077.28 449549.89
\end{verbatim}

Verification against the lm() function The manually computed
coefficients are virtually identical to those obtained with
\texttt{lm()} in Exercise 2. This confirms that our implementation
correctly reproduces the OLS estimators.

Modification of function We extend our OLS function by adding the
boolean argument \texttt{add\_constant}, which allows us to estimate the
regression either with or without an intercept.

\begin{Shaded}
\begin{Highlighting}[]
    \CommentTok{\#In this case "false" {-}\textgreater{} in order to verify beta\_1 with the exercise 3 above}
\NormalTok{estimate\_ols\_modified }\OtherTok{\textless{}{-}} \ControlFlowTok{function}\NormalTok{(y, x, }\AttributeTok{add\_constant =} \ConstantTok{FALSE}\NormalTok{) \{ }
  
  \ControlFlowTok{if}\NormalTok{ (add\_constant }\SpecialCharTok{==} \ConstantTok{TRUE}\NormalTok{) \{}
    
    \CommentTok{\# OLS with intercept}
\NormalTok{    beta\_1 }\OtherTok{\textless{}{-}} \FunctionTok{sum}\NormalTok{((x }\SpecialCharTok{{-}} \FunctionTok{mean}\NormalTok{(x)) }\SpecialCharTok{*}\NormalTok{ (y }\SpecialCharTok{{-}} \FunctionTok{mean}\NormalTok{(y))) }\SpecialCharTok{/} \FunctionTok{sum}\NormalTok{((x }\SpecialCharTok{{-}} \FunctionTok{mean}\NormalTok{(x))}\SpecialCharTok{\^{}}\DecValTok{2}\NormalTok{)}
\NormalTok{    beta\_0 }\OtherTok{\textless{}{-}} \FunctionTok{mean}\NormalTok{(y) }\SpecialCharTok{{-}}\NormalTok{ beta\_1 }\SpecialCharTok{*} \FunctionTok{mean}\NormalTok{(x)}
    
    \FunctionTok{return}\NormalTok{(}\FunctionTok{c}\NormalTok{(}\AttributeTok{beta\_0 =}\NormalTok{ beta\_0, }\AttributeTok{beta\_1 =}\NormalTok{ beta\_1))}
    
\NormalTok{  \} }\ControlFlowTok{else}\NormalTok{ \{}
    
    \CommentTok{\# OLS without intercept}
\NormalTok{    beta\_1 }\OtherTok{\textless{}{-}} \FunctionTok{sum}\NormalTok{(x }\SpecialCharTok{*}\NormalTok{ y) }\SpecialCharTok{/} \FunctionTok{sum}\NormalTok{(x}\SpecialCharTok{\^{}}\DecValTok{2}\NormalTok{)}
    
    \FunctionTok{return}\NormalTok{(beta\_1)}
\NormalTok{  \}}
\NormalTok{\}}

\FunctionTok{estimate\_ols\_modified}\NormalTok{(df}\SpecialCharTok{$}\NormalTok{gdp, df}\SpecialCharTok{$}\NormalTok{surname)}
\end{Highlighting}
\end{Shaded}

\begin{verbatim}
## [1] 532247.2
\end{verbatim}

Verification with result in subexercise 4 The estimated coefficient
obtained with \texttt{add\_constant\ =\ FALSE} is identical to the
result from
\texttt{lm(gdp\ \textasciitilde{}\ surname\ -\ 1,\ data\ =\ df)} in
subexercise 4 (532247), confirming that the implementation correctly
reproduces the OLS estimator for the model without an intercept.

\end{document}
